\chapter{Complexity — Alg Circuit Example}
%%% generated by the blueprint pipeline from TCSlib.Complexity.Randomized.AlgCircuitExample
\section{Overview}

Declarations of the Lean module \texttt{TCSlib.Complexity.Randomized.AlgCircuitExample}.

\section{Declarations}

\begin{definition}[Arithmetic-operation gates]
\lean{Randomized.AlgGate.isOp}
\statementsource{Vad12}{pdf-bf2ce5422088-p016-b001}
\label{Randomized.AlgGate.isOp}
\leanok
\uses{Randomized.AlgGate}
For a gate in a circuit with $n$ input variables, the Boolean-valued indicator is true
exactly for addition, subtraction, and multiplication gates, and false for
input-variable and constant gates.
\end{definition}

\begin{definition}[Variable-or-one gate]
\lean{Randomized.AlgCircuit.varOrOne}
\label{Randomized.AlgCircuit.varOrOne}
\leanok
\uses{Randomized.AlgGate}
For a circuit with $n$ input variables and a natural number $i$, choose the
input-variable gate indexed by $i$ if $i<n$; otherwise choose the constant-one gate.
\end{definition}

\begin{definition}[Product circuit gate list]
\lean{Randomized.AlgCircuit.prodOnePlusGates}
\label{Randomized.AlgCircuit.prodOnePlusGates}
\leanok
\uses{Randomized.AlgCircuit.varOrOne, Randomized.AlgGate}
For each natural number $k$, define a gate list recursively. At $k=0$ the list consists
of the constant-one gate; at $k+1$, append to the list for $k$ the gate for input $x_k$
when $k<n$ (and the constant-one gate otherwise), followed by an addition gate with node
indices $0$ and $3k+1$, and a multiplication gate with node indices $3k$ and $3k+2$.
\end{definition}

\begin{definition}[Product circuit]
\lean{Randomized.AlgCircuit.prodOnePlus}
\label{Randomized.AlgCircuit.prodOnePlus}
\leanok
\uses{Randomized.AlgCircuit, Randomized.AlgCircuit.prodOnePlusGates}
For each natural number $n$, the product circuit with $n$ input variables is the circuit
whose ordered gate list is the recursively defined product gate list of length parameter
$n$.
\end{definition}

\begin{theorem}[Length of the product circuit]
\lean{Randomized.AlgCircuit.length_prodOnePlusGates}
\label{Randomized.AlgCircuit.length_prodOnePlusGates}
\leanok
\uses{Randomized.AlgCircuit.prodOnePlusGates}
\difficulty{3}
For any natural number $k$, the recursively constructed gate list has length exactly
$3k+1$.
\end{theorem}

\begin{theorem}[Operation count in the generated gate list]
\lean{Randomized.AlgCircuit.countP_isOp_prodOnePlusGates}
\label{Randomized.AlgCircuit.countP_isOp_prodOnePlusGates}
\leanok
\uses{Randomized.AlgCircuit.prodOnePlusGates, Randomized.AlgCircuit.varOrOne, Randomized.AlgGate.isOp}
\difficulty{4}
For every natural number $n$ and every natural number $k$, the gate list constructed for
$k$ factors contains exactly $2k$ arithmetic operation gates, counting addition,
subtraction, and multiplication gates.
\end{theorem}

\begin{theorem}[Earlier references in the constructed gate list]
\lean{Randomized.AlgCircuit.refsBelow_prodOnePlusGates}
\label{Randomized.AlgCircuit.refsBelow_prodOnePlusGates}
\leanok
\uses{Randomized.AlgCircuit.length_prodOnePlusGates, Randomized.AlgCircuit.prodOnePlusGates, Randomized.AlgCircuit.varOrOne, Randomized.AlgGate.refsBelow}
\difficulty{6}
For a natural number $n$, let $G_k$ be the gate list defined by $G_0=[1]$ and by
appending, at each step from $k$ to $k+1$, the gates $	ext{varOrOne}(k)$, addition of
gates $0$ and $3k+1$, and multiplication of gates $3k$ and $3k+2$. For every pair of
natural numbers $k,i$ with $i<|G_k|$, the gate at position $i$ is either an input
variable or the constant $1$, or is an addition, subtraction, or multiplication gate
whose two referenced positions are both strictly less than $i$.
\end{theorem}

\begin{theorem}[Well-formedness of the product circuit]
\lean{Randomized.AlgCircuit.wellFormed_prodOnePlus}
\label{Randomized.AlgCircuit.wellFormed_prodOnePlus}
\leanok
\uses{Randomized.AlgCircuit.WellFormed, Randomized.AlgCircuit.length_prodOnePlusGates, Randomized.AlgCircuit.prodOnePlus, Randomized.AlgCircuit.refsBelow_prodOnePlusGates}
\difficulty{3}
For every nonnegative integer $n$, the circuit constructed as the product of one plus
each of its inputs is well formed: it has at least one gate, and every gate that refers
to earlier nodes refers only to nodes with smaller indices than its own position.
\end{theorem}

\begin{theorem}[Operation count in the product circuit]
\lean{Randomized.AlgCircuit.countP_isOp_prodOnePlus}
\label{Randomized.AlgCircuit.countP_isOp_prodOnePlus}
\leanok
\uses{Randomized.AlgCircuit.countP_isOp_prodOnePlusGates, Randomized.AlgCircuit.prodOnePlus, Randomized.AlgGate.isOp}
\difficulty{1}
For every natural number $n$, the constructed circuit for $\prod_{i=1}^{n}(1+X_i)$
contains exactly $2n$ arithmetic operation gates, counting addition, subtraction, and
multiplication gates.
\end{theorem}

\begin{theorem}[Size of the product circuit]
\lean{Randomized.AlgCircuit.size_prodOnePlus}
\label{Randomized.AlgCircuit.size_prodOnePlus}
\leanok
\uses{Randomized.AlgCircuit.length_prodOnePlusGates, Randomized.AlgCircuit.prodOnePlus, Randomized.AlgCircuit.size}
\difficulty{1}
For every natural number $n$, the circuit constructed to compute the product of the
first $n$ factors has size $3n+1$.
\end{theorem}

\begin{theorem}[Lookup at the appended value]
\lean{Randomized.AlgCircuit.getD_append_singleton_length}
\label{Randomized.AlgCircuit.getD_append_singleton_length}
\leanok
\difficulty{1}
Let $R$ be a type with a specified zero, let $vs$ be a finite list of elements of $R$,
and let $v$ be an element of $R$. Looking up index equal to the length of $vs$ in the
list obtained by appending $v$ to $vs$, with $0$ as the default value, returns $v$.
\end{theorem}

\begin{theorem}[Evaluation of a product circuit]
\lean{Randomized.AlgCircuit.run_prodOnePlusGates}
\label{Randomized.AlgCircuit.run_prodOnePlusGates}
\leanok
\uses{Randomized.AlgCircuit.getD_append_singleton_length, Randomized.AlgCircuit.length_prodOnePlusGates, Randomized.AlgCircuit.length_run, Randomized.AlgCircuit.prodOnePlusGates, Randomized.AlgCircuit.run, Randomized.AlgCircuit.run_append_singleton, Randomized.AlgCircuit.varOrOne, Randomized.AlgGate, Randomized.AlgGate.eval}
\difficulty{5}
Let $R$ be a commutative ring, let $x$ assign an element of $R$ to each of $n$ inputs,
and let $k$ be a natural number. In the successive evaluation of the circuit built to
compute the product of the first $k$ factors, the value at index $0$ (taken to be $0$ if
that index is absent) is $1$, and the value at index $3k$ (again taken to be $0$ if
absent) is $\prod_{i=0}^{k-1}(1+v_i)$, where $v_i=x_i$ if $i<n$ and $v_i=1$ otherwise.
\end{theorem}

\begin{theorem}[Polynomial of the product circuit]
\lean{Randomized.AlgCircuit.toMvPolynomial_prodOnePlus}
\label{Randomized.AlgCircuit.toMvPolynomial_prodOnePlus}
\leanok
\uses{Randomized.AlgCircuit.eval, Randomized.AlgCircuit.prodOnePlus, Randomized.AlgCircuit.run_prodOnePlusGates, Randomized.AlgCircuit.size_prodOnePlus, Randomized.AlgCircuit.toMvPolynomial, Randomized.AlgCircuit.varOrOne, Randomized.AlgGate, Randomized.AlgGate.eval}
\difficulty{3}
For every natural number $n$, the integer-coefficient multivariate polynomial
represented by the circuit constructed to compute the product of the first $n$ factors
is $\prod_{i\in\mathrm{Fin}(n)}(1+X_i)$, where $X_i$ is the variable corresponding to
input $i$.
\end{theorem}

\begin{theorem}[Support size of a product polynomial]
\lean{Randomized.AlgCircuit.card_support_prod_one_add_X}
\label{Randomized.AlgCircuit.card_support_prod_one_add_X}
\leanok
\difficulty{6}
For every natural number $n$, the product $\prod_{i=1}^{n}(1+X_i)$ in the integer
polynomial ring in variables $X_1,\ldots,X_n$ has exactly $2^n$ monomials with nonzero
coefficient.
\end{theorem}

\begin{theorem}[Support size of a circuit polynomial]
\lean{Randomized.AlgCircuit.card_support_toMvPolynomial_prodOnePlus}
\label{Randomized.AlgCircuit.card_support_toMvPolynomial_prodOnePlus}
\leanok
\uses{Randomized.AlgCircuit.card_support_prod_one_add_X, Randomized.AlgCircuit.prodOnePlus, Randomized.AlgCircuit.toMvPolynomial, Randomized.AlgCircuit.toMvPolynomial_prodOnePlus}
\difficulty{1}
For every natural number $n$, the polynomial computed by the constructed arithmetic
circuit for $\prod_{i=1}^{n}(1+X_i)$ has exactly $2^n$ monomials with nonzero
coefficient.
\end{theorem}
